
%\section{Introduction}\label{sec:intro}

Microlocal Morse theory is machinery 
to study the propagation of cohomology 
of complexes of sheaves. 
More generally, this theory is considered as 
a part of microlocal sheaf theory, 
due to Kashiwara and Schapira. 
They defined the notion of 
microsupport of a complex of sheaves 
in the 1980's. 

To explain more precisely, 
let us fix some notations. 
We consider a commutative unital ring \(\kk\) 
of finite global dimension (e.g., \(\zz\) or a field). 
Let \(X\) be a topological space. 
We denote by \(\Dompb(\kk_X)\) (resp.\ \(\Dompl(\kk_{X})\)) 
the bounded derived category (resp.\ bounded 
derived category from below) of 
sheaves of \(\kk\)-modules on \(X\).

If \(M\) is a real manifold of class \(C^{\infty}\) 
and \(F\) a complex of sheaves, 
then the microsupport \(\muS[](F)\) of \(F\) is 
a closed conic (i.e., stable 
under \(\rr^{+}\)-action) involutive (i.e., coisotropic) subset 
of the cotangent bundle of \(M\). 
The precise definition is given in Section \ref{sec:t-ss}, 
but roughly speaking, \(\muS[](F)\) describes the (co)directions 
of \(X\) in which the cohomology of \(F\) cannot be extended. 
This notion can be interpreted by analogy with Morse theory. 
Let \(\varphi\) be a real-valued function of class \(C^1\). 
Then we can replace 
the notion of critical points of \(\varphi\) with 
the one with respect to \(F\). 
That is to say, we can define a critical point 
with respect to \(F\) by the point \(p\) 
such that \(d\varphi(p)\) belongs to \(\muS[](F)\). 
We can then obtain an analogous result 
of Morse theory in a
framework of sheaf theory. 
In Section \ref{sec:t-morse}, we prove a truncated version 
of (a particular case of) 
the microlocal Morse lemma (cf.\ \cite[Proposition 1.8]{GKS12} 
or \cite[Corollary 5.4.19]{KS90}). 
Roughly speaking, this asserts that 
if there is no critical point with respect 
to a complex of sheaves, 
we have the natural isomorphism of the cohomologies 
of the sublevel sets of the manifold.

%\subsection{Truncated Microsupport}}

%However, we replace the notion of microsupports with
In 2003, 
Kashiwara, Monteiro-Fernandes, and Schapira
introduced the notion of truncated microsupport in \cite{KMS03a}, 
in the course of the study of solutions of 
boundary value problems of \(\mcal{D}\)-modules. 
The difference between 
microsupports and truncated microsupports 
is as follows. 
On the one hand, 
a point \(p\) in the cotangent bundle \(\cotb[M]\) does 
not belong to the microsupport if and only if 
a certain kind of cohomologies of 
all the degrees vanish. 
On the other hand, the point \(p\) does 
not belong to the truncated microsupport if and only if 
the same kind of cohomologies of all the degrees less than 
a fixed number vanish.

After they introduced truncated microsupports, 
some properties 
of truncated microsupports were shown. 
For example, 
a truncated microsupport 
is not involutive, but
has a property called weak involutivity. 
This is shown also 
by Kashiwara, Monteiro-Fernandes, 
and Schapira in \cite{KMS03b}. 
Some functorial properties of truncated microsupport are 
studied by several authors. See \cite{KMS03a,FM07}. 
These authors studied truncated microsupports
mainly in the context of analytic problems.

%\subsection{Results}}

In this paper, 
in contrast to the previous research, 
we shall apply the truncated microsupport 
to the study of some geometrical problems. 
First we prove a truncated version of 
the noncharacteristic deformation lemma. 
Consider a Hausdorff space and a family of open subsets of it. 
Assume that the family is ``noncharacteristic'' with 
respect to a complex of sheaves on the space.
In this case, 
the original version of the theorem asserts that 
the global section on a member of the family 
is isomorphic to the one on the union of the family. 
Note that this isomorphism is one between derived functors.

%
This lemma is one of the most fundamental 
results in microlocal sheaf theory, and 
used throughout \cite{KS90}.
Thus a truncated version of this lemma 
is also the key result in this paper. 
However, one should notice that 
the original proof of the lemma cannot be applied 
to our truncated version. 
To prove the result, 
we make three modifications 
to the lemma. 
First, we add an assumption on a total space. 
That is, 
we assume that the space is not only Hausdorff, but also 
satisfies the condition that 
all the open subspaces are paracompact.
Second, we weaken the hypothesis of vanishing of 
cohomologies. In the original lemma, we assume that 
the cohomologies of all degrees vanish. However, in our version, 
we assume that the cohomologies of not all degrees 
but degrees less than a certain number vanish.
Third, we make an additional assumption on 
the family of deformation.
Therefore, 
to make, in particular, the third
assumption acceptable to the reader, 
we give several examples of families of deformation with figures. 
Then, in turn, we formulate 
a truncated version (Theorem \ref{thm:t-morse}) of
the microlocal Morse lemma (\cite[Proposition 5.4.17]{KS90}), 
under the light of truncated microsupport. 
Note that we can also study 
the unique continuation property of 
cohomology groups of sheaves 
through the truncated microsupport.
%Note that by using the truncated microsupport, we can gain the 
%information of the unique continuation. 
More precisely, if \(k\) is an integer and 
we consider the truncated microsupport \(\muS[k](F)\) of 
a bounded complex of sheaves \(F\), 
we can investigate whether 
the cohomology sheaves \(H^{j}(F)\) propagate or not for \(j<k\).
%in particular, 
Furthermore, we can see 
whether the cohomology sheaves \(H^k(F)\) of the \(k\)-th degree 
enjoy the unique continuation property or not.


%whether the continuation \(H^{k}(F)\) is unique or not, 
%if it can be extended. 

%In \cite{KMS03a}, they do not study the functorial property of 
%truncated microsupport of direct images of non proper morphisms. 
%We prove this functorial property under non proper direct images. 

%\subsection{Organization}
%\subsection{Organization}}

The paper is organized as follows. 
In Section \ref{sec:review}, 
we review the notations and results from \cite{KS90}. 
In Section \ref{sec:t-ss}, 
we define truncated microsupport of sheaves, 
due to \cite{KMS03a}. 
In Section \ref{sec:t-nonchar}, 
we prove a truncated version of 
the noncharacteristic deformation lemma. 
In the proof, Lemma \ref{clm:ind-nonchar} plays a crucial role. 
However, since the proof of the lemma is long 
and rather technical, 
its proof will be given in Section \ref{sec:proof-lem}.
Finally in Section \ref{sec:t-morse}, 
we apply the result in Section \ref{sec:t-nonchar} 
to the microlocal Morse theory. 
%In Section \ref{sec:t-nonchar}, 
%we prove Lemma \ref{clm:ind-nonchar}, which is too long to 
%write in the proof of Theorem \ref{prp:t-nonchar}. 
